% Figura: retratos de fase de 1/z, 1/z^2 y 1/(z^2 (z-1))
% Se usa sombreado plano (una celda = un color) para no interpolar a través
% del salto de +180 a -180 grados del argumento.
\begin{figure}[ht]
\centering
\begin{tikzpicture}
  \begin{groupplot}[group style={group size=3 by 1, horizontal sep=0.55cm},
      oscuro, grid=none, view={0}{90}, width=5.15cm, height=5.15cm,
      axis equal image, colormap name=fase, shader=flat corner,
      point meta min=-180, point meta max=180,
      xmin=-1.5, xmax=1.5, ymin=-1.5, ymax=1.5,
      xtick={-1,0,1}, ytick={-1,0,1},
      samples=\resF, samples y=\resF, domain=-1.5:1.5, y domain=-1.5:1.5,
      tick label style={font=\tiny}, title style={yshift=-2pt},
      xlabel={$\Real z$}, xlabel style={font=\scriptsize, yshift=4pt}]
    % (a) 1/z: arg = -atan2(y,x)
    \nextgroupplot[title={(a) $1/z$\,: polo simple},
        ylabel={$\Imag z$}, ylabel style={font=\scriptsize, yshift=-4pt}]
      \addplot3[surf, point meta={-atan2(y,x)}] (x, y, 0);
      \addplot3[only marks, mark=x, mark size=3pt, line width=1.2pt, blanco]
        coordinates {(0,0,0)};
    % (b) 1/z^2: arg = -2 atan2(y,x), reducido a (-180,180]
    \nextgroupplot[title={(b) $1/z^{2}$\,: polo doble}, yticklabels={}]
      \addplot3[surf, point meta={mod(-2*atan2(y,x)+540,360)-180}] (x, y, 0);
      \addplot3[only marks, mark=x, mark size=3pt, line width=1.2pt, blanco]
        coordinates {(0,0,0)};
    % (c) 1/(z^2 (z-1)): arg = -2 atan2(y,x) - atan2(y,x-1)
    \nextgroupplot[title={(c) $1/\bigl(z^{2}(z-1)\bigr)$}, yticklabels={},
        colorbar right, colorbar style={width=0.22cm, ytick={-180,0,180},
          yticklabels={$-\pi$,$0$,$\pi$}, tick label style={font=\tiny, text=gris},
          axis line style={grismedio}, ylabel={$\arg f$},
          ylabel style={font=\scriptsize, text=texto, yshift=-2pt}}]
      \addplot3[surf, point meta={mod(-2*atan2(y,x)-atan2(y,x-1)+900,360)-180}] (x, y, 0);
      \addplot3[only marks, mark=x, mark size=3pt, line width=1.2pt, blanco]
        coordinates {(0,0,0) (1,0,0)};
  \end{groupplot}
\end{tikzpicture}
\caption{Retratos de fase: cada punto $z$ se colorea según $\arg f(z)$
(barra de la derecha). Alrededor de un polo de orden $p$ (marcado con
$\times$) el círculo cromático se recorre $p$ veces: una vez en (a), dos en
(b). En (c) conviven un polo doble en $z=0$ y uno simple en $z=1$, que son
los del \cref{ej:polos}\,(ii).}
\label{fig:fases}
\end{figure}
